The stress harness ran \lvl{low} and \lvl{max} on seven families at $n=20$, $100$ and $400$
(\val{stressRuns} runs, Table~\ref{tab:stress}). Every tour was checked to be a valid permutation with
a correctly recomputed length. No run produced an invalid tour (\val{stressInvalid} failures), and no
tour came out shorter than the proven bound (\val{stressViolations} violations).

\begin{table}[t]
\centering\footnotesize
\caption{Degenerate and known-optimum families. ``Gap to bound'' is listed per $n$ ($20$, $100$, $400$)
and is a ceiling on the true gap. \emph{fine} and \emph{coarse} are the same point set scaled by
$10^{-3}$ and $10^{6}$.}
\label{tab:stress}
\setlength{\tabcolsep}{4pt}
\begin{tabular}{@{}llrrrl@{}}
\toprule
family & level & $n$ & runs & invalid & result\\
\midrule
convex & \lvl{low} & 20, 100, 400 & 3 & 0 & hull cycle (optimal) at every $n$\\
convex & \lvl{max} & 20, 100, 400 & 3 & 0 & hull cycle (optimal) at every $n$\\
line & \lvl{low} & 20, 100, 400 & 3 & 0 & known optimum at every $n$\\
line & \lvl{max} & 20, 100, 400 & 3 & 0 & known optimum at every $n$\\
grid & \lvl{low} & 20, 100, 400 & 3 & 0 & gap to bound 0.00\%, 0.00\%, +0.41\%\\
grid & \lvl{max} & 20, 100, 400 & 3 & 0 & gap to bound 0.00\%, 0.00\%, 0.00\%\\
dup & \lvl{low} & 20, 100, 400 & 3 & 0 & gap to bound 0.00\%, +1.21\%, +2.32\%\\
dup & \lvl{max} & 20, 100, 400 & 3 & 0 & gap to bound 0.00\%, +0.50\%, +0.85\%\\
same & \lvl{low} & 20, 100, 400 & 3 & 0 & known optimum at every $n$\\
same & \lvl{max} & 20, 100, 400 & 3 & 0 & known optimum at every $n$\\
fine & \lvl{low} & 20, 100, 400 & 3 & 0 & gap to bound 0.00\%, +0.95\%, +3.84\%\\
fine & \lvl{max} & 20, 100, 400 & 3 & 0 & gap to bound 0.00\%, +0.94\%, +1.65\%\\
coarse & \lvl{low} & 20, 100, 400 & 3 & 0 & gap to bound 0.00\%, +0.95\%, +3.96\%\\
coarse & \lvl{max} & 20, 100, 400 & 3 & 0 & gap to bound 0.00\%, +0.94\%, +1.65\%\\
\bottomrule
\end{tabular}

\end{table}

\paragraph{Known optima by construction.}
For points in convex position the hull cycle is optimal, and both levels returned it at every size.
At $n=1000$, \lvl{low} returned \val{convex1000@low@hull} in \val{convex1000@low} and \lvl{max}
returned \val{convex1000@max@hull} in \val{convex1000@max} of CPU time. Collinear points (optimum: twice
the span) and coincident points (optimum: 0) were solved exactly. On the $20\times 20$ grid,
\lvl{max} met the 1-tree bound, which proves its tour optimal, while \lvl{low} did not.

\paragraph{The bound can be very loose.}
On collinear points the tour is optimal by construction, yet its bound gap reads
\val{stress@line@max@100} at $n=100$ and \val{stress@line@max@400} at $n=400$. A degree-constrained relaxation cannot express the structure of a line. This
is the same phenomenon as fl417 in Section~\ref{sec:tsplib}, in a more extreme form.

\paragraph{Scale dependence.}
The \emph{fine} and \emph{coarse} families are one point set at two scales, so a scale-free method
would return the same relative gap on both. \lvl{max} nearly does (\val{stress@fine@max@400} against \val{stress@coarse@max@400} at
$n=400$). \lvl{low} does not (\val{stress@fine@low@400} against \val{stress@coarse@low@400}), because the anneal temperature is set in absolute units (Section~\ref{sec:limits}).

\paragraph{Determinism.}
Every level was run three times on each of two $n=100$ instances (\val{repeatRuns} runs in
\val{repeatGroups} groups). In \val{repeatIdentical} of the \val{repeatGroups} groups the repeats
returned the same length and the same tour. Runs are a deterministic function of (instance, seed,
level), also under the worker pool.
